\section{Method}
\label{sec:method}

\hh{} keeps Hydra's pipeline (dilated convolutions on the series and its first difference,
competing groups, counts of winners, sparse scaling and a ridge classifier) and changes
only the groups. Each group is the vertex set of a Hanner polytope, given by a labelled
cotree on the kernel's taps. This section defines the groups
(Section~\ref{sec:groups}), shows that the winner of a group follows the cotree
(Section~\ref{sec:recursion}), interprets the resulting features
(Section~\ref{sec:interpretation}), and describes the sampler, the transform, the
operation count, the controls and the implementation
(Sections~\ref{sec:sampling}--\ref{sec:implementation}).

\subsection{Hanner groups}
\label{sec:groups}

A kernel has $\ell$ taps, and a patch is a vector $x\in\R^\ell$ of $\ell$ dilated samples of
the input. A \emph{labelled cotree} $T$ on the taps is a full binary tree whose leaves are
the taps $1,\dots,\ell$, each exactly once, and whose internal nodes are labelled $\lone$
(an $\ell_1$ sum) or $\linf$ (an $\ell_\infty$ product). Every subtree $S$ covers a set
$\tau(S)$ of taps, and the two children of a node cover disjoint sets.

\begin{definition}[Hanner group]
\label{def:group}
The \emph{group} of a subtree $S$ is the ordered list $V(S)$ of vectors in $\R^\ell$
defined by
\begin{itemize}
\item a leaf on tap $i$: $V(S)=(-e_i,\,e_i)$;
\item $S=S_1\lone S_2$: the list $V(S_1)$ followed by $V(S_2)$;
\item $S=S_1\linf S_2$: the vectors $v_a+u_b$ for $v_a\in V(S_1)$, $u_b\in V(S_2)$,
  ordered by the index $a\,n_2+b$, where $n_2=|V(S_2)|$ and indices start at 0.
\end{itemize}
The group of the cotree is $V=V(T)$, and its size is $|V|$.
\end{definition}

The convex hull of $V$ is the Hanner polytope of $T$: a leaf gives the segment
$[-e_i,e_i]$, a $\linf$ node gives the Cartesian product $K_1\times K_2$ of the
children's polytopes (they live in complementary coordinate subspaces), and a $\lone$ node
gives $\operatorname{conv}(K_1\cup K_2)$. These are the two operations that define Hanner
polytopes \citep{hanner1956intersections,hansen1981structure}.

\begin{proposition}[Structure of a group]
\label{prop:structure}
For every labelled cotree $T$ on $\ell$ taps:
\begin{enumerate}
\item the vectors in $V$ are distinct, and they are exactly the vertices of
  $\operatorname{conv}V$;
\item every vector has entries in $\{-1,0,1\}$, and its support is a union of whole
  leaves;
\item $V=-V$: the map $\alpha$ defined by swapping the two vertices of each leaf, acting
  on each part of a $\lone$ node and on both factors of a $\linf$ node, is an involution
  with $V_{\alpha(j)}=-V_j$;
\item $2\ell\le|V|\le 2^\ell$; the cross-polytope (all nodes $\lone$) attains the lower
  bound and the cube (all nodes $\linf$) the upper one.
\end{enumerate}
\end{proposition}

\begin{proof}
All four statements follow by induction over the tree. (1)~The vertices of a product of
polytopes are the pairs of vertices. For two polytopes in complementary subspaces that each
contain the origin in their relative interior, the vertices of the convex hull of their union
are the vertices of the two polytopes. Distinctness follows because the parts of a $\lone$
node have disjoint supports and the parts of a $\linf$ node are recovered by restricting to
$\tau(S_1)$ and $\tau(S_2)$. (2)~Each tap belongs to one leaf, and a $\linf$ node adds vectors
with disjoint supports. (3)~Negation commutes with both operations. (4)~A subtree on $t$
taps has between $2t$ and $2^t$ vertices: a leaf has 2; a $\lone$ node has $n_1+n_2$ and a
$\linf$ node $n_1n_2$ vertices, and $n_1+n_2\le n_1n_2$ whenever $n_1,n_2\ge2$. The bounds
are not characterisations: a product of two leaves has $4=2+2$ vertices, so, for example,
$(x_1\linf x_2)\lone x_3$ also has $2\ell=6$ vertices.
\end{proof}

\subsection{Support-function recursion}
\label{sec:recursion}

Hydra's group responds to a patch $x$ with the winning kernel
$\operatorname{arg\,max}_j\langle x,w_j\rangle$ and its response. For a Hanner group, the
response of the winner is the support function $h(x)=\max_{v\in V}\langle x,v\rangle$ of
the polytope, and it follows the cotree.

\begin{proposition}[Winner recursion]
\label{prop:recursion}
Compute, bottom-up over the cotree, a value $h(S)$ and an index $w(S)$ for every subtree:
\begin{itemize}
\item leaf on tap $i$: $h=|x_i|$; $w=1$ (the vertex $e_i$) if $x_i>0$, and $w=0$ (the
  vertex $-e_i$) otherwise;
\item $S=S_1\linf S_2$: $h=h(S_1)+h(S_2)$ and $w=w(S_1)\,n_2+w(S_2)$;
\item $S=S_1\lone S_2$: if $h(S_1)\ge h(S_2)$, then $h=h(S_1)$ and $w=w(S_1)$; otherwise
  $h=h(S_2)$ and $w=n_1+w(S_2)$, with $n_1=|V(S_1)|$.
\end{itemize}
At the root, $h(T)=\max_{v\in V}\langle x,v\rangle$, and $w(T)$ is the lowest index in $V$
at which this maximum is attained. The recursion uses $\ell$ absolute values and $\ell-1$
additions or comparisons of real numbers, and no multiplications; each leaf's sign test
$x_i>0$ yields both $|x_i|$ (by a select) and the leaf's winner.
\end{proposition}

\begin{proof}
By induction, with the claim that $h(S)=\max_{v\in V(S)}\langle x,v\rangle$ and that $w(S)$
is the lowest maximising index in $V(S)$.
\emph{Leaf.} The responses are $-x_i$ (index 0) and $x_i$ (index 1). If $x_i>0$, only
index~1 attains $|x_i|$; if $x_i<0$, only index~0; if $x_i=0$, both do and the lowest is~0.
\emph{Product.} Since $v_a$ and $u_b$ have disjoint supports,
$\langle x,v_a+u_b\rangle=\langle x,v_a\rangle+\langle x,u_b\rangle$, so the maximum is
$h(S_1)+h(S_2)$ and the set of maximisers is the product of the children's maximiser sets.
Because $0\le b<n_2$, the index $a\,n_2+b$ is lexicographic in $(a,b)$, so its minimum
over a product set is attained at the two smallest indices.
\emph{Sum.} $V(S)$ is the union of the two lists, with the first list at indices below
$n_1$, so the maximum is $\max(h(S_1),h(S_2))$. If $h(S_1)\ge h(S_2)$, the first list
contains a maximiser, and every index in it is lower than every index in the second.
Otherwise every maximiser is in the second list, at offset $n_1$.
\emph{Count.} A binary tree with $\ell$ leaves has $\ell-1$ internal nodes, each costing one
addition or one comparison. The index updates are integer arithmetic on the winners and do
not involve the data.
\end{proof}

The index $w(T)$ is the first maximiser in list order, which is what a brute-force argmax
over $V$ returns, so the recursion and a brute-force evaluation give the same winner and
the same counts, including ties. This
statement is about exact arithmetic. In floating point, the recursion adds the responses in
the order of the tree, while a dot product adds them in tap order. Our reference
implementation and the fast kernels both use the recursion's order and agree bit for bit
(Section~\ref{sec:implementation}).

\begin{proposition}[Minimum and hard-min counts]
\label{prop:hardmin}
For a Hanner group, $\min_{v\in V}\langle x,v\rangle=-h(x)$, attained at $\alpha(w)$. If the
minimiser is taken to be $\alpha(w)$, the hard-min count of every vertex $v$ equals the
hard-max count of $-v$, so the hard-min features are a fixed permutation of the hard-max
features. The same holds for every centrally symmetric group.
\end{proposition}

This follows from $V=-V$ (Proposition~\ref{prop:structure}). The tie rule matters: Hydra
takes the \emph{first} minimiser in list order, which differs from $\alpha(w)$ whenever the
minimum is attained more than once. Ties are not rare: a patch with a zero tap, as at the
zero-padded ends of a series, ties at that leaf. With Hydra's rule the hard-min counts are
therefore a permutation of the hard-max counts only up to tied patches. Our implementation
uses $\alpha(w)$. A per-feature scaler and a linear classifier are invariant to a fixed
permutation of the features, so these hard-min counts carry no information beyond the
hard-max counts, and \hh{} counts soft maxima and hard maxima (Section~\ref{sec:transform}).
The ablation with hard-min counts (Section~\ref{sec:results}) uses the same $\alpha(w)$
rule; its identical predictions confirm the permutation argument in the implementation,
but they are not an independent test of Hydra's first-minimiser rule.

\subsection{What the features compute}
\label{sec:interpretation}

The two extreme cotrees give two familiar codes. For the cube ($|V|=2^\ell$, all
$\linf$), the winner is the sign pattern of the patch,
$(\operatorname{sign}x_1,\dots,\operatorname{sign}x_\ell)$ with $0$ mapped to $-1$, and $h$
is $\|x\|_1$. This is the code of hyperplane hashing \citep{charikar2002simhash}, here with
the coordinate axes as hyperplanes. For the cross-polytope ($|V|=2\ell$, all $\lone$), the
winner is the coordinate of largest magnitude with its sign, and $h=\|x\|_\infty$. This is
the code of cross-polytope hashing \citep{andoni2015crosspolytope}, without its random
rotation. In a mixed cotree, a $\linf$ node concatenates the winners of its children and a
$\lone$ node keeps the winner of the child with the larger response. The winner is
therefore a sign pattern on a subset of the taps, and the $\lone$ nodes choose the subset
by where the patch responds most strongly. In \hh{} the randomness that hashing gets from
random projections comes from the sampled cotree, the dilation and the channel subset.

The features are therefore close to counts of binary symbolic words, as in dictionary
methods \citep{lin2007sax,schaefer2023weasel2}, and to sign patterns of neighbouring
samples \citep{kaya2014lbp}. They differ from ordinal patterns \citep{bandt2002permutation},
which record the rank order of the samples rather than their signs. The soft-max count
weights each word by the response $h$, which, for a cube on centred taps, is the patch's
$\ell_1$ deviation from its mean.

\subsection{Cotree sampling}
\label{sec:sampling}

Groups are drawn at random under a vertex budget $V_{\max}$. The sampler starts from the
$\ell$ leaves in random order and merges two randomly chosen subtrees until one tree
remains. A merge of subtrees with $n_a$ and $n_b$ vertices is a $\linf$ product with
probability $p_{\text{product}}$ if
\[
  R + n_a n_b \le V_{\max},
\]
where $R$ is the total vertex count of the other subtrees; otherwise it is a $\lone$ sum.
Merging all remaining subtrees by sums would give $R+n_an_b$ vertices, the smallest
possible final count, so the condition keeps the final group within $V_{\max}$ and no draw
is ever rejected. Every tap appears in every group, so $2\ell\le|V|\le V_{\max}$. With the
frozen setting ($\ell=9$, $V_{\max}=32$, $p_{\text{product}}=0.5$), groups have 18 to 32
vertices, 24.4 on average (Table~\ref{tab:operations}). The sampler can also restrict
merges to adjacent blocks of taps; this option is off in all reported runs.

\subsection{The transform}
\label{sec:transform}

The transform follows Hydra \citep{dempster2023hydra} in everything but the groups.
\begin{itemize}
\item \emph{Dilations and inputs.} The dilations are $2^0,2^1,\dots,
  2^{\lfloor\log_2((T-1)/(\ell-1))\rfloor}$ for series of length $T$, with zero padding
  that keeps the output length. Each dilation is applied to the series and to its first
  difference.
\item \emph{Groups per block.} For each pair of dilation and input (a block), groups are
  drawn until their vertices, the \emph{atoms}, reach a budget of 256, the number of
  kernels Hydra uses per block ($k=8$ kernels in 32 groups). A block has 10.9 groups on
  average.
\item \emph{Channels.} For a series with $C>1$ channels, each group draws
  $\operatorname{clip}(\lfloor C/2\rfloor,2,8)$ channels at random (with replacement), the
  subset size used by Hydra, and works on their sum.
\item \emph{Centring.} In each block, every second group sees mean-centred patches
  $x-\bar{x}\mathbf{1}$, and the others see raw patches. The centred groups respond to the
  shape of a patch regardless of its level; the raw groups also see the level relative to
  zero (ablation in Section~\ref{sec:results}).
\item \emph{Features.} For every vertex $v$ of every group, the transform records the
  number of time steps at which $v$ wins (hard-max count) and the sum of the winning
  responses $h$ at those steps (soft-max count). A block contributes two features per
  atom.
\item \emph{Classifier.} The features go through Hydra's sparse scaler, which takes square
  roots, standardises each feature and sets the entries that were zero back to zero. A
  ridge classifier follows, with its regularisation strength chosen by leave-one-out
  cross-validation among 10 values from $10^{-3}$ to $10^{3}$.
\end{itemize}

\subsection{Operation count}
\label{sec:operations}

\begin{table}[t]
\centering\small
% Generated by experiments/paper_tables.py; do not edit by hand.
\begin{tabular}{lrrrrr}
\toprule
 & Vertices & Operations & Groups & Operations & Multipli- \\
 & per group & per group & per block & per block & cations \\
\midrule
Hydra (Gaussian) & 8 & 150 & 32 & 4,800 & 2,304 \\
Matched control, as timed (dense) & 24.4 & 438.8 & 10.9 & 4,804 & 2,408 \\
Matched control, sparse (not timed) & 24.4 & 54.4 & 10.9 & 596 & 0 \\
Hanner-Hydra (recursion) & 24.4 & 17 & 10.9 & 186 & 0 \\
\midrule
\multicolumn{6}{l}{Ratios per block: Hydra / Hanner 25.8$\times$, dense control / Hanner 25.8$\times$, sparse control / Hanner 3.2$\times$} \\
\multicolumn{6}{l}{(2,000 sampled blocks, seed 0)} \\
\bottomrule
\end{tabular}

\caption{Arithmetic per time step of one block (one dilation and input, 256 atoms), under
the frozen sampler. Operations are additions, subtractions, comparisons and, for Hydra,
multiplications; \hh{} also takes $\ell=9$ absolute values per group, included in its
count. The matched control has the same vertex sparsity as \hh{}. As timed, it evaluates
every vertex as a dense dot product over all $\ell$ taps, multiplying by $-1$, $0$ or $1$;
the sparse row counts an evaluation over the non-zeros only, which we did not implement.
Hydra finds both the maximum and the minimum of each group. Patch extraction, channel
sums, centring and counting are shared by all arms and not counted.}
\label{tab:operations}
\end{table}

Table~\ref{tab:operations} counts the arithmetic that finding the winners needs: the
intrinsic cost of a group, not a model of the measured times. A Hydra group of 8 Gaussian
kernels of length 9 costs 72 multiplications, 64 additions and 14 comparisons for 8 atoms.
A Hanner group costs 17 operations for 24 atoms on average, with no multiplication. A
ternary group with the same vertices costs 439 operations when every vertex is evaluated
densely, as in our control kernel, and 54 when only the non-zeros are visited. Per block,
Hydra and the dense control both need about $26\times$ the operations of \hh{}, and a
sparse control would need $3.2\times$.

The measured transform-time ratios do not follow these counts: Hydra / \hh{} is
$11.2\times$, \texttt{control\_sym} / \hh{} is $3.8\times$, and Hydra / \texttt{control\_sym}
is $3.0\times$, although Hydra and the dense control have nearly the same count
(Section~\ref{sec:results}). The shared costs in the caption (patch extraction, channel
sums and counting, which touch memory) take a large share of the time, and Hydra runs its
convolutions in PyTorch, in aeon's implementation. Two consequences for the comparison with
the controls: the $3.8\times$ compares the recursion with a dense ternary evaluation of the
same groups, so it is an upper bound on what the cotree adds over the best ternary
implementation; and the counts alone do not say how much of it a sparse control kernel would
recover. We did not time one.

\emph{What ``multiplication-free'' covers.} It covers the group evaluation: the winner and
its response need only absolute values, additions and comparisons. The first difference is
a subtraction, and soft-max counts add the winning responses. Centring is the one place
where our implementation divides: it subtracts the patch mean $\bar{x}=s/\ell$, with
$s=\sum_i x_i$, once per patch, not once per group or vertex. A multiplication-free form
exists: evaluate the group on $\ell x_i-s$ instead, where $9x_i=8x_i+x_i$ is a shift and an
addition for $\ell=9$. Since $\langle\ell x-s\mathbf{1},v\rangle=\ell\,\langle
x-\bar{x}\mathbf{1},v\rangle$, the winners and hard-max counts are unchanged in exact
arithmetic, and the soft-max counts of the centred groups are scaled by $\ell$, which the
per-feature scaler removes. The sparse scaler and the ridge classifier do multiply, as in
Hydra and MiniRocket. The kernels we timed use floating point and the division; all
timings are for that implementation, not for the scaled form.

\subsection{Theory-matched controls}
\label{sec:controls}

Two controls separate the cotree structure from the other properties of a Hanner group.
For every Hanner group, a control draws a random group of ternary vectors with the same
size and the same number of non-zeros per vertex: the positions and signs of the non-zeros
are uniform, and duplicate vectors are redrawn. \texttt{control} draws the vectors
independently. \texttt{control\_sym} draws them in antipodal pairs $(u,-u)$ placed at the
pair positions of the Hanner group, so it is centrally symmetric and has the same antipode
map $\alpha$. Both controls are evaluated by ternary dot products, $|V|$ per patch, each over all $\ell$
taps (Section~\ref{sec:operations}), and
pass through the same transform, scaler and classifier. A difference between \hh{} and
\texttt{control\_sym} can therefore come only from the cotree structure, and a difference
between the two controls only from central symmetry. The pre-registration names
\texttt{control\_sym} as the primary control (Section~\ref{sec:protocol}).

\subsection{Implementation}
\label{sec:implementation}

The reference implementation is written in NumPy. It fixes the order of every
floating-point sum, so that a second implementation can reproduce it exactly. The timed
implementation consists of single-threaded Numba kernels, one for Hanner groups (the
recursion of Proposition~\ref{prop:recursion}) and one for the ternary controls. The tests
check that both kernels reproduce the reference bit for bit, and that the recursion agrees
with brute force over the vertex list, including ties. The method is also available as
aeon estimators \citep{middlehurst2024aeon}, a transformer and a classifier that wrap the
reference and reproduce it exactly. Section~\ref{sec:reproducibility} lists the code and
data.
